\lead{Synthesis.} Fine-tuning has two regimes. With few gentle passes or a frozen body the pretrained related-DNA features transfer and give the best numbers in this report (1.228 after one pass at LR $3\times10^{-4}$; 1.243 with the body frozen for 16 passes), where every scratch recipe stops at 1.259 to 1.263. With the scratch recipe (LR .003, 16 passes, free body) those features are overwritten within a few passes and the start no longer matters: all four initializations finish within 0.004 of one another. Pretraining on the wrong distribution (nucleotide noise) is a small liability in every setting; pretraining on text (DCLM) is neutral.

\lead{Rules of thumb.} From a related pretrained model, protect the body: freeze it, pull toward it, or use a $10\times$ lower LR for one to four passes; never reuse the scratch LR with 16 passes, and never run many passes at WD 0. Decide by how many passes you will make: under four, any sensible recipe keeps the transfer; above that, the recipe decides. Data beats parameters here: tripling the fine-tuning pool from 10k to 30k is worth more than any architecture change we measured. Keep WD at the tuned value regardless of initialization.
